\subsection{The Syntax of Regular Expressions}
\label{sec:regex-syntax}
\slidesrc{Syntax: RE \& word: 3--8, 10}

\begin{nbdef}{Regular expression}
A \textbf{regular expression} (RE, regex) is a formal language for specifying text
strings. A regex is a \emph{pattern}; searching a text with it returns the substrings
that the pattern \emph{matches}.
\end{nbdef}

The motivating question: how can we search for any of \emph{woodchuck},
\emph{woodchucks}, \emph{Woodchuck}, \emph{Woodchucks} with a single query? A plain string
search for ``woodchuck'' misses the capitalised forms; we need a way to say ``either
\texttt{w} or \texttt{W}'' and ``an optional \texttt{s}''.

\subsubsection{Disjunction: square brackets and the pipe}
\slidesrc{Syntax: RE \& word: 4, 6}

Letters inside square brackets mean \textbf{one character among these}:

\begin{center}
\small
\begin{tabular}{@{}l l l@{}}
\toprule
\textbf{Pattern} & \textbf{Matches} & \textbf{Example} \\
\midrule
\re{[wW]oodchuck}   & \emph{Woodchuck} or \emph{woodchuck} & \\
\re{[1234567890]}   & any digit & \\
\re{[A-Z]}          & an upper-case letter (a \textbf{range}) & \underline{D}renched Blossoms \\
\re{[a-z]}          & a lower-case letter & \underline{m}y beans were impatient \\
\re{[0-9]}          & a single digit & Chapter \underline{1}: Down the Rabbit Hole \\
\bottomrule
\end{tabular}
\end{center}

Brackets choose among single \emph{characters}. To choose among longer strings use the
\textbf{pipe} \re{|}:

\begin{center}
\small
\begin{tabular}{@{}l l@{}}
\toprule
\textbf{Pattern} & \textbf{Matches} \\
\midrule
\re{groundhog|woodchuck} & \emph{groundhog} or \emph{woodchuck} (they are the same animal!) \\
\re{yours|mine} & \emph{yours} or \emph{mine} \\
\re{a|b|c} & the same as \re{[abc]} \\
\re{[gG]roundhog|[Ww]oodchuck} & both animals, capitalised or not \\
\bottomrule
\end{tabular}
\end{center}

\subsubsection{Negation inside brackets}
\slidesrc{Syntax: RE \& word: 5}

A caret \re{\textasciicircum} right after the opening bracket negates the set: \re{[\textasciicircum
Ss]} means ``any character that is \emph{neither} \texttt{S} nor \texttt{s}''.

\begin{center}
\small
\begin{tabular}{@{}l l l@{}}
\toprule
\textbf{Pattern} & \textbf{Matches} & \textbf{Example} \\
\midrule
\re{[\textasciicircum A-Z]} & not an upper-case letter & O\underline{y}fn pripetchik \\
\re{[\textasciicircum Ss]}  & neither \texttt{S} nor \texttt{s} & \underline{I} have no exquisite reason \\
\re{[\textasciicircum e\textasciicircum]} & neither \texttt{e} nor \texttt{\textasciicircum} & \underline{L}ook here \\
\re{a\textasciicircum b}    & the literal string \texttt{a\textasciicircum b} & Look up \underline{a\textasciicircum b} now \\
\bottomrule
\end{tabular}
\end{center}

\begin{nbwarning}[the caret has three meanings]
\re{\textasciicircum} means \emph{negation} only when it is the \emph{first} character inside
\re{[\,]}. In \re{[e\textasciicircum]} it is an ordinary character, in \re{a\textasciicircum b}
it is an ordinary character, and at the start of a pattern (\re{\textasciicircum[A-Z]}) it
is an \emph{anchor} (see below).
\end{nbwarning}

\subsubsection{Repetition: \texttt{?}, \texttt{*}, \texttt{+} and the wildcard \texttt{.}}
\slidesrc{Syntax: RE \& word: 7}

These operators apply to the character (or group) right before them:

\begin{center}
\small
\begin{tabular}{@{}l l l@{}}
\toprule
\textbf{Pattern} & \textbf{Meaning} & \textbf{Matches} \\
\midrule
\re{colou?r} & previous character optional (0 or 1) & \emph{color}, \emph{colour} \\
\re{oo*h!}   & 0 or more of the previous character & \emph{oh!}, \emph{ooh!}, \emph{oooh!}, \ldots \\
\re{o+h!}    & 1 or more of the previous character & \emph{oh!}, \emph{ooh!}, \emph{oooh!}, \ldots \\
\re{baa+}    & \texttt{ba} followed by one or more \texttt{a} & \emph{baa}, \emph{baaa}, \emph{baaaa}, \ldots \\
\re{beg.n}   & \texttt{.} is any single character & \emph{begin}, \emph{begun}, \emph{beg3n} \\
\bottomrule
\end{tabular}
\end{center}

The star \re{*} and the plus \re{+} are called the \textbf{Kleene star} and
\textbf{Kleene plus}, after the mathematician Stephen C.~Kleene who introduced regular
expressions. Note that \re{oo*h!} and \re{o+h!} describe exactly the same set of strings:
$x^{+}$ is just $x x^{*}$.

\subsubsection{Anchors: \texttt{\textasciicircum} and \texttt{\$}}
\slidesrc{Syntax: RE \& word: 8}

Anchors do not match a character but a \emph{position}: \re{\textasciicircum} the start
of the line, \re{\$} the end.

\begin{center}
\small
\begin{tabular}{@{}l l l@{}}
\toprule
\textbf{Pattern} & \textbf{Matches} & \textbf{Example} \\
\midrule
\re{\textasciicircum[A-Z]} & an upper-case letter at the start of a line & \underline{P}alo Alto \\
\re{\textasciicircum[\textasciicircum A-Za-z]} & a non-letter at the start of a line & \underline{1}\quad \underline{``}Hello'' \\
\re{\textbackslash.\$} & a full stop at the end of a line & The end\underline{.} \\
\re{.\$} & any character at the end of a line & The end\underline{?}\quad The end\underline{!} \\
\bottomrule
\end{tabular}
\end{center}

In \re{\textbackslash.\$} the backslash \textbf{escapes} the dot: \re{\textbackslash.}
is a literal full stop, while \re{.} alone is the wildcard.

\subsubsection{Summary of the operators}
\slidesrc{Syntax: RE \& word: 10}

\begin{center}
\small
\begin{tabularx}{\linewidth}{@{}l l l X@{}}
\toprule
\textbf{Operator} & \textbf{Meaning} & \textbf{Example} & \textbf{Example meaning} \\
\midrule
\re{+}     & one or more   & \re{a+}    & 1 or more \texttt{a} characters \\
\re{*}     & zero or more  & \re{a*}    & 0 or more \texttt{a} characters \\
\re{?}     & optional      & \re{a?}    & 0 or 1 \texttt{a} characters \\
\re{[\,]}  & choose 1      & \re{[abc]} & \texttt{a} or \texttt{b} or \texttt{c} \\
\re{[-]}   & range         & \re{[a-z]} & any character between \texttt{a} and \texttt{z} \\
\re{[\textasciicircum]} & not & \re{[\textasciicircum a]} & any character that is not \texttt{a} \\
\re{(\,)}  & grouping      & \re{([a-z])+} & one or more characters between \texttt{a} and \texttt{z} \\
\re{(|)}   & or            & \re{(ey|ax)} & the string \texttt{ey} or the string \texttt{ax} \\
\re{ab}    & follow        & \re{ab}    & \texttt{a} followed by \texttt{b} \\
\re{\textasciicircum} & start & \re{\textasciicircum a} & \texttt{a} at the start of the string/line \\
\re{\$}    & end           & \re{a\$}   & \texttt{a} at the end of the string/line \\
\re{\textbackslash s} & whitespace & \re{\textbackslash sa} & a whitespace followed by \texttt{a} \\
\re{.}     & any character & \re{a.b}   & \texttt{a}, any character, \texttt{b} \\
\bottomrule
\end{tabularx}
\end{center}

\begin{nbwarning}[a typo in the summary slide]
The slide gives \re{(a-z)+} as the grouping example. Inside round brackets the hyphen is
\emph{not} a range: \re{(a-z)+} matches the literal three-character string
\texttt{a-z} (repeated). The range needs square brackets: \re{([a-z])+} or simply
\re{[a-z]+}.
\end{nbwarning}

Tools to experiment: \url{https://regex101.com/} shows step by step what a pattern
matches and explains each part of it.
